\chapter{INTRODUCTION}

\section{Motivation}

Large language models (LLMs) deliver strong performance on reasoning, coding, and generation tasks, but they remain expensive to serve, slower to scale, and demanding in terms of hardware and energy consumption. Recent work therefore explores whether small language models (SLMs) can be combined with larger models through routing, planning, debate, retrieval, or verification so that a system can retain most of the quality of an LLM while using fewer expensive calls in practice \cite{hao2024hybrid, shao2025dot, oh2025uncertainty, multiagent2025}. This thesis project is motivated by that systems question rather than by a single-model leaderboard question.

The central problem is not simply ``which model is best?'' but rather ``which architecture is best under a shared workload and a shared measurement policy?'' In a real deployment setting, accuracy alone is not sufficient. Two systems with similar answer quality may differ substantially in latency, token usage, cost, power draw, carbon impact, and operational fragility. This gap is especially important for coding tasks, where users expect both correctness and responsiveness. The separate systematic literature review completed within this project showed that the research community has produced many collaborative SLM--LLM designs, yet the reporting of cost and energy remains much weaker than the reporting of raw task accuracy \cite{slrmanuscript2026, energytoken2025}. That observation directly motivates the design of the current platform.

\section{Project Scope and Research Goal}

The thesis project implements an experiment platform in which every architecture receives the same \texttt{Query} object and returns the same \texttt{Response} object. This common contract makes it possible to compare standalone, hybrid, ensemble, and swarm-style strategies under the same benchmark conditions. The project therefore treats architecture as the primary independent variable. Models still matter, but they are studied as components inside a larger orchestration surface.

The current report focuses on four tightly connected goals:

\begin{itemize}
    \item building a reproducible experiment platform for SLM/LLM comparison,
    \item grounding the design in a completed systematic literature review,
    \item defining a coding-focused benchmark and evaluation protocol for the next phase of experimentation, and
    \item introducing project-specific architecture ideas and measurement choices, including the EATS metric.
\end{itemize}

Training and fine-tuning are intentionally outside the scope of this report. Although the repository contains training-related material, this term report does not treat training as part of the completed experiment pipeline and does not use it to support any claim.

\section{Contributions of the Current Report}

The contribution of this report is not a final empirical result set. Instead, it is a design-complete and implementation-aware draft of the thesis work at its current stage. Concretely, the report contributes the following:

\begin{itemize}
    \item a documented experiment platform that integrates remote model serving, a local control plane, experiment execution, reporting, and a web-based management interface;
    \item a measurement methodology that treats latency, token usage, cost, energy, and carbon as first-class outcomes rather than secondary diagnostics;
    \item a coding-benchmark plan centered on a project-specific \texttt{custom\_stratified} dataset with easy, medium, and hard difficulty levels; and
    \item a clear separation between implemented platform components, preliminary experimental components, and project-proposed swarm architectures that still require validation.
\end{itemize}

\section{Structure of the Report}

Chapter 2 summarizes the relevant literature and compresses the findings of the completed systematic literature review into a term-report-friendly narrative. Chapter 3 presents the methodology, including the software stack, the hardware topology, the architecture taxonomy, the benchmark semantics, and the measurement design. Chapter 4 defines the experimental study as a design and execution plan rather than as a final result chapter, because the current experiments are not yet valid for thesis-grade claims. Chapter 5 concludes with the current status of the project and the next milestones required to reach a valid comparative evaluation.
